\section{Discussion}

\subsection{Five Clinical Mysteries Resolved}
The IEC framework provides mechanistic answers to five long-standing clinical mysteries regarding adolescent idiopathic scoliosis (AIS):
1) \emph{Why does AIS emerge during adolescence?} The Energy Deficit Window reveals that the metabolic cost of maintaining high-anisotropy mechanosensors (demand) scales steeply ($L^4$) during rapid growth, outstripping the diffusion-limited supply ($L^2$).
2) \emph{Why do scoliotic curves typically progress in a specific order?} The VIM Cascade predicts a sequential failure of structural proteins—Vimentin fails first (at 80\% supply), followed by Lamin A/C, then PIEZO2, blinding the spine to its own curvature.
3) \emph{Why are lateral curves (scoliosis) more common than purely sagittal deformities?} The Vector-Scalar Mismatch mechanism explains that under predominantly axial loading, lateral perturbations drive stronger misalignment between the developmental information gradient and the stress vector, collapsing the local effective stiffness.
4) \emph{Why do girls progress to severe curves 10 times more frequently than boys?} Our metabolic modeling suggests girls experience an earlier and deeper metabolic depth deficit ($R_{peak} \approx 2.7$ vs $2.4$ in boys), trapping the system in a positive-feedback loop of sensor failure and asymmetric growth.
5) \emph{Why does the S-curve persist in astronauts?} By framing curvature as an information-selected minimum rather than a passive gravitational response, our framework naturally predicts that the underlying genetic target shape persists in microgravity, albeit with degraded stiffness due to convective shutdown.

\subsection{IEC versus Competing Theories of AIS}
Three theories currently compete to explain AIS: the estrogen-growth hypothesis~\cite{weinstein2008adolescent}, the melatonin hypothesis, and the proprioceptive mismatch hypothesis~\cite{assaraf2020piezo2}.
The estrogen hypothesis explains the sex disparity but not curve patterns or why only some rapidly growing girls develop scoliosis. The melatonin hypothesis (melatonin deficiency leading to bone growth asymmetry) was not confirmed in clinical trials. The proprioceptive hypothesis explains why PIEZO2 mutations cause scoliosis but does not explain the timing of onset during adolescence.

The IEC framework subsumes all three. Estrogen timing determines the duration of the high-velocity growth phase; melatonin modulates ARNTL/BMAL1 circadian coupling that serves as a critical supply factor; and PIEZO2 failure is explicitly predicted as a downstream consequence in the VIM Cascade. IEC thus provides a unifying thermodynamic envelope explaining when and why all existing mechanisms converge to produce macroscopic deformity.

\subsection{The Mechanism: Energy Deficits and Vector Mismatch}
The protein-level foundation of IEC involves two key insights. First, mechanosensory and structural proteins compete for metabolic resources, creating energy deficit windows during rapid growth where mechanical feedback is transiently compromised. Second, information-elasticity coupling is fundamentally vectorial: tissue resistance to deformation depends on the alignment between developmental gradients and stress directions. These mechanisms explain why scoliosis emerges during adolescence (energy deficit timing) and why deviations are predominantly lateral (vector mismatch in that direction).

\subsection{Testable predictions and experimental validation}
Our framework makes twelve falsifiable predictions, spanning molecular genetics, environmental physiology, and clinical biomechanics. We highlight the most critical:

1. \textbf{Hoxd11 conditional knockout in lumbar somites}: Lumbar lordosis decreases from wild-type $50 \pm 5^\circ$ to $30 \pm 5^\circ$ (mouse P21 micro-CT).
2. \textbf{$\chi_\kappa$-stratified scoliosis progression}: Patients with model-inferred $\chi_\kappa > 0.08$ m$^{-1}$ will progress $>2\times$ faster than those with $\chi_\kappa < 0.05$ m$^{-1}$ over 2 years (n=200 cohort).
3. \textbf{ISS astronaut lordosis}: Serial MRI should show $<20\%$ reduction in lumbar lordosis after 6 months in orbit, versus passive beam predictions of $>80\%$ loss.
4. \textbf{Spinal clock phase shift}: Urinary sulfatoxymelatonin phase in AIS patients will be shifted $>2$ hours from healthy controls and correlate with Cobb angle ($r>0.4$).
5. \textbf{Mode-pattern prediction from muscle asymmetry}: Preoperative paraspinal MRI quantification of regional LBX1/VIM ratio asymmetry should predict Lenke curve type with AUC $> 0.75$.

\subsection{Priority Experiments}
To validate the core claims of this framework, three priority experiments are required:
1. \textbf{PIEZO2 proprioceptive neuron CKO during adolescent growth (mouse, 6 weeks)}: Measure VIM $\to$ LMNA $\to$ LBX1 expression cascade in paraspinal muscle at weekly intervals. \emph{Predicted result}: VIM network fragmentation appears 2 weeks before Cobb angle development.
2. \textbf{Paraspinal muscle biopsies from AIS patients at peak height velocity}: Measure PPARGC1A/PGC-1$\alpha$, PIEZO2 density, and LMNA nuclear aspect ratio. \emph{Predicted result}: Concave paraspinal PPARGC1A expression will be $\geq 20\%$ lower than the convex side in actively progressing patients.
3. \textbf{Asymmetric skin tension in axis elongation}: Unilateral collagenase application to dorsal chick skin during axis elongation should induce spine curvature toward the low-tension side, demonstrating information gradient misalignment.

\subsection{Limitations and Model Assumptions}
Our framework is purely theoretical and computational. No wet-laboratory experiments were conducted in this study. The protein anisotropy values derive from AlphaFold predictions, which may differ from experimentally resolved structures in complex in vivo environments. The scaling exponents ($L^2$, $L^4$) are theoretically motivated by macroscopic growth assumptions but not directly measured at the cellular level. The energy deficit parameters ($R_{peak} \approx 2.7$ in females vs $2.4$ in males) are model estimates, not clinical measurements.

Furthermore, our current implementation treats protein dynamics with simplified kinetics and does not fully capture complex mechanochemical feedback loops. Direct experimental measurement of mechanosensory protein turnover rates during growth spurts would significantly refine these predictions. Finally, our tensor coupling model assumes orthotropic symmetry; future work must incorporate full anisotropic elasticity tensors measured via multi-directional mechanical testing. All predictions require prospective experimental validation before any clinical application.
